\documentclass[11pt]{article}
\usepackage[a4paper,margin=27mm]{geometry}
\usepackage[T1]{fontenc}
\usepackage{lmodern,amsmath,amssymb,amsthm,booktabs,array,longtable}
\usepackage[colorlinks=true,linkcolor=blue,citecolor=blue,urlcolor=blue]{hyperref}
\newtheorem{theorem}{Theorem}
\newtheorem{lemma}[theorem]{Lemma}
\newtheorem{corollary}[theorem]{Corollary}
\newcommand{\rev}{\operatorname{rev}}
\newcommand{\pl}{\operatorname{pl}}
\newcommand{\cl}{\operatorname{cpl}}
\title{A signed-weight extension of Frid's recurrence for the Thue--Morse prefix problems of Ambro\v{z}, Kadlec, Mas\'akov\'a and Pelantov\'a}
\author{Haobo Ma\\\small ChronoAI Pte Ltd, Singapore\and
Wenlin Zhang\\\small National University of Singapore\\\small ChronoAI Pte Ltd, Singapore}
\date{18 November 2026}
\begin{document}
\maketitle
\begin{abstract}
Extending Frid's ordinary-length recurrence (2019, Theorem~4), we prove
the sharp factor-two and unbounded equal-cost assertions of Ambro\v{z},
Kadlec, Mas\'akov\'a and Pelantov\'a, arXiv:1812.00711v2, Section~6,
p.~13, final paragraph and Tables~1--2, p.~14, by identifying the combined
pal-antipal length of each zero-start Thue--Morse prefix with the minimum
weight of a signed binary expansion of its length.
Desubstitution proves this formula and the sharp comparison with ordinary
palindromic length. We determine all first equal-cost occurrences and the
exact logarithmic upper growth rate of the combined length.
Frid's recurrence also yields prefixes attaining the factor-two endpoint
at every positive combined cost.
\end{abstract}
\noindent\small This note is licensed under Creative Commons Attribution 4.0
International (CC BY 4.0).\normalsize

\section{The prefix question and the result}
Let $t=t_0t_1\cdots$ be the binary Thue--Morse word, with
\[
 t_0=0,\qquad t_{2m}=t_m,\qquad t_{2m+1}=1-t_m\quad(m\geq0).
\]
For a binary word $w$, let $E(w)$ exchange $0$ and $1$ letter by letter.
A palindrome satisfies $w=\rev(w)$; an antipalindrome satisfies
$w=E(\rev(w))$. Define $\pl(w)$ and $\cl(w)$ as the minimum numbers of
consecutive nonempty factors in a factorization of $w$, allowing respectively
palindromes, and palindromes or antipalindromes. Both costs of the empty word
are zero. Put
\[
 P(n)=\pl(t_0\cdots t_{n-1}),\qquad
 C(n)=\cl(t_0\cdots t_{n-1})\quad(n\geq0).
\]
These are costs of the prefix starting at zero, not maxima over all factors.
For an integer $z$, let $W(z)$ be the least $k\geq0$ such that
\[
 z=\sum_{h=1}^{k}\varepsilon_h2^{e_h},\qquad
 \varepsilon_h\in\{-1,1\},\quad e_h\in\mathbb Z_{\geq0}.
\]
Repeated exponents are allowed in this definition.

The final paragraph of Section~6 of \cite{AKMP} suggests that allowing
antipalindromes reduces the prefix factor count ``by half (and not more)''.
It then proposes arbitrarily large equal ordinary and combined costs.
These are unnumbered suggestions in a prefix context, rather than assertions
about all shifted factors. We prove both, including a first-occurrence
formula for the equal-cost prefixes displayed in Table~2.

\begin{theorem}\label{thm:main}
For every $n\geq0$,
\[
 C(n)=W(n),\qquad C(n)\leq P(n)\leq2C(n).
\]
Both endpoints are attained at every positive combined cost: for every
$k\geq1$ there are positive prefix lengths $a,b$ with
$P(a)=C(a)=k$ and $P(b)=2C(b)=2k$.
\end{theorem}

Frid \cite[Lemma~3 and Theorem~4]{Frid}, first posted in June 2019,
determines $P$ by recurrences modulo four. Li \cite[Introduction and
Proposition~1]{Li}, posted in July 2019, gives later independent
ordinary-length results and generalizations, acknowledging Frid's work.
The signed-weight sequence itself is classical: see OEIS A007302
\cite{OEIS}, including Wilson's 2005 and Fox's 2013 recurrences, and
\cite[Section~2, p.~3]{KW} on minimum-weight nonadjacent forms.
Our contribution is its identification with the mixed minimum $C$.
The palindrome/antipalindrome interchange and centered desubstitution
already occur in \cite[Section~5, Proposition~6 and Lemmas~2--3]{BBGL}.
The proof below uses the two
halvings of an odd-start even-length factor; no recurrence for $P$ is needed
for that formula or for the universal factor-two bound. We use Frid's exact
recurrence only to give sharp examples at every combined cost.

\section{Signed weights and one-letter changes}
\begin{lemma}\label{lem:weight}
We have $W(0)=0$ and $W(1)=1$. For $z\in\mathbb Z$ and $m\geq0$,
\begin{align}
 W(-z)&=W(z),
 &|W(z+1)-W(z)|&\leq1,\label{eq:basic}\\
 W(2m)&=W(m),
 &W(2m+1)&=1+\min\{W(m),W(m+1)\}.\label{eq:binary}
\end{align}
Consequently,
\begin{align}
 W(4m)&=W(m),&W(4m+1)&=W(m)+1,\nonumber\\
 W(4m+2)&=1+\min\{W(m),W(m+1)\},
 &W(4m+3)&=W(m+1)+1.\label{eq:four}
\end{align}
\end{lemma}
\begin{proof}
Changing all signs gives symmetry. Concatenating representations proves
$W(x+y)\leq W(x)+W(y)$, which gives \eqref{eq:basic}.
In a minimum representation, two opposite unit terms cancel and two
equal unit terms combine into one term of exponent one. Either operation
would reduce the number of terms. Thus there are no unit terms when the
integer is even, and exactly one, of sign $1$ or $-1$, when it is odd.
Dividing the remaining terms by two gives the lower bounds in
\eqref{eq:binary}; doubling a minimum representation and, in the odd case,
adding a unit term gives the reverse inequalities.
For example,
$W(4m+1)=1+\min\{W(m),1+\min(W(m),W(m+1))\}$.
The second entry of the outer minimum is at least $W(m)$ by
\eqref{eq:basic}. This proves the residue-one formula. The residue-three
formula follows in the same way with $m+1$ in place of $m$; the even
residues follow directly from \eqref{eq:binary}.
\end{proof}

For either cost $L=\pl$ or $L=\cl$, appending a letter changes the cost
by at most one:
\begin{equation}\label{eq:lip}
 |L(wa)-L(w)|\leq1.
\end{equation}
Appending a singleton proves one direction. For the other, take an optimal
factorization of $wa$. If its last factor is a singleton, remove it.
Otherwise write that last factor as $bva$. Removing its two outside
letters leaves a palindrome or antipalindrome $v$ of the same type.
After deleting the final $a$, replace $bva$ by the singleton $b$ followed
by $v$, omitting $v$ if empty. This increases the number of factors by
at most one. In the ordinary case all factors involved are palindromes.

\section{A lower bound on every mixed cut}
Call a nonempty factor allowed if it is a palindrome or an antipalindrome.
For half-open intervals write $t[i:j)=t_i\cdots t_{j-1}$.

\begin{lemma}\label{lem:cut}
If $0\leq i<j$ and $t[i:j)$ is allowed, then
\begin{equation}\label{eq:cut}
 W(j)\leq W(i)+1.
\end{equation}
\end{lemma}
\begin{proof}
Let $\mu(0)=01$ and $\mu(1)=10$. Then $t=\mu(t)$ and
\[
 \rev(\mu(v))=\mu(E(\rev(v))).
\]
Since $\mu$ is injective on binary words and commutes with $E$,
$\mu(v)$ is allowed if and only if $v$ is allowed, with palindrome and
antipalindrome interchanged.

There are no constant factors of length three: every three consecutive
positions contain an aligned pair $t_{2a},t_{2a+1}$ of different letters.
There is no palindrome of length five either. If such a factor starts
at $2a$, its five letters are
\[
 t_a,\ 1-t_a,\ t_{a+1},\ 1-t_{a+1},\ t_{a+2};
\]
the palindrome equations force $t_a=t_{a+1}=t_{a+2}$.
Starting at $2a+1$ gives
$1-t_a,t_{a+1},1-t_{a+1},t_{a+2},1-t_{a+2}$ and the same contradiction.
Every odd palindrome of length at least five contains a central
palindrome of length five. Antipalindromes have even length. Hence an
allowed odd-length factor has length one or three.

We now use strong induction on $j-i$. Lengths one and two satisfy
\eqref{eq:cut} by subadditivity and $W(1)=W(2)=1$.
For length three the factor is a palindrome. Its first and third letters
cannot agree if $i\equiv0$ or $1\pmod4$, by the defining recurrence for $t$.
If $i=4a+2$, \eqref{eq:four} gives
\[
 W(i)=1+\min\{W(a),W(a+1)\},\qquad W(j)=1+W(a+1),
\]
so \eqref{eq:basic} gives \eqref{eq:cut}.
If $i=4a+3$, the values are $1+W(a+1)$ and
$1+\min\{W(a+1),W(a+2)\}$, so the increment is nonpositive.

It remains to consider even length $j-i=2k\geq4$.
If $i=2a$, put $j=2b$. Then $t[i:j)=\mu(t[a:b))$,
and the induction hypothesis at the shorter allowed factor $t[a:b)$,
together with $W(2a)=W(a)$ and $W(2b)=W(b)$, gives the result.

If $i=2a+1$ and $j=2b+1$, then $b-a=k$.
Removing the first and last letters leaves
$\mu(t[a+1:b))$, an allowed word. Extending by one letter on each side
also gives an allowed word, namely $\mu(t[a:b+1))$.
To check extension, write the original reflection equation as
$1-t_a=E^\varepsilon(t_b)$, where $\varepsilon\in\{0,1\}$ and
$E^0$ is the identity on letters. It implies
$t_a=E^\varepsilon(1-t_b)$, the equation for the two added letters;
all inner reflection equations are unchanged.
The two desubstituted words have lengths $k-1$ and $k+1$, both positive
and smaller than $2k$. Thus induction gives
\[
 W(b)\leq W(a+1)+1,\qquad W(b+1)\leq W(a)+1.
\]
It follows that
\[
 \min\{W(b),W(b+1)\}\leq\min\{W(a),W(a+1)\}+1.
\]
Apply the odd formula in \eqref{eq:binary} to $i$ and $j$ to finish.
\end{proof}

For any mixed factorization of a prefix, let
$0=i_0<i_1<\cdots<i_q=n$ be its cuts. Summing
$W(i_h)-W(i_{h-1})\leq1$ gives $W(n)\leq q$.
Taking the minimum proves
\begin{equation}\label{eq:lower}
 W(n)\leq C(n).
\end{equation}

\section{The exact mixed minimum and the comparison}
Applying $\mu$ to a mixed factorization preserves the number of allowed
nonempty factors, so $C(2m)\leq C(m)$. By \eqref{eq:lip},
\[
 C(2m+1)\leq1+\min\{C(m),C(m+1)\}.
\]
Strong induction, starting with $C(0)=0$ and $C(1)=1$, and
\eqref{eq:binary} prove $C(n)\leq W(n)$. For odd $n\geq3$ both
indices $m,m+1$ are smaller than $n$; the initial odd case is $n=1$.
Together with \eqref{eq:lower} this proves $C(n)=W(n)$.

Put $\tau=\mu^2$, so $\tau(0)=0110$ and $\tau(1)=1001$.
Both letter images are palindromes, and reversal commutes with $\tau$.
Consequently $\tau$ carries a palindromic factorization to one with
the same number of factors, and $P(4m)\leq P(m)$.
Using \eqref{eq:lip} from the neighboring multiples of four gives
\begin{align*}
 P(4m+1)&\leq P(m)+1,\\
 P(4m+2)&\leq\min\{P(m),P(m+1)\}+2,\\
 P(4m+3)&\leq P(m+1)+1.
\end{align*}
Strong induction and \eqref{eq:four} now give $P(n)\leq2W(n)$.
The empty case is immediate. In residue two the upper bound is
$2\min\{W(m),W(m+1)\}+2=2W(4m+2)$; in residues one and three
the added one is at most the added two in $2W$.
Finally $C\leq P$ because every ordinary factorization is mixed.
The prefix $01$ has $C(2)=1$ and $P(2)=2$, so the constant two is optimal.

\section{Attainment at both endpoints}
\begin{corollary}\label{cor:equal}
Set $a_0=1$ and $a_{r+1}=4a_r-1$. Then for every $r\geq0$,
\[
 a_r=\frac{2\cdot4^r+1}{3},\qquad P(a_r)=C(a_r)=r+1.
\]
\end{corollary}
\begin{proof}
Deleting one letter from the prefix of length $4a_r$ gives
$P(a_{r+1})\leq P(a_r)+1$, so $P(a_r)\leq r+1$.
The residue-three formula gives
$W(4a_r-1)=W(a_r)+1$. Thus $C(a_r)=W(a_r)=r+1$,
and $C\leq P$ gives equality. The closed formula follows from the
recurrence and $a_0=1$.
\end{proof}

\begin{corollary}\label{cor:first}
For $k\geq1$, put $g_k=(2\cdot4^{k-1}+1)/3$. Then
\[
 \min\{n\geq1:C(n)=k\}
 =\min\{n\geq1:P(n)=C(n)=k\}=g_k.
\]
Thus every first occurrence in Table~2 of \cite{AKMP} follows from the
same formula, not just its finitely many displayed rows.
\end{corollary}
\begin{proof}
Corollary~\ref{cor:equal} gives $P(g_k)=C(g_k)=k$.
We prove $W(n)\leq k-1$ for $0\leq n<g_k$ by induction on $k$.
For $k=1$ this concerns only $n=0$.
Since $g_{k+1}=4g_k-1$, write $n<g_{k+1}$ as $4m+r$,
where $0\leq r\leq3$. Then $m<g_k$; if $r=3$, also $m+1<g_k$.
In \eqref{eq:four}, residues zero and one use $W(m)\leq k-1$,
residue two uses $\min\{W(m),W(m+1)\}\leq W(m)$,
and residue three uses $W(m+1)\leq k-1$.
Hence $W(n)\leq k$ in every case. Now use $C=W$.
The numbers $g_k$ are OEIS A007583, with index $k-1$; its 2023
comment describes their signed-weight extremal meaning \cite{OEIS}.
The conclusion here is their equal-cost prefix interpretation.
\end{proof}

\begin{corollary}\label{cor:growth}
For the zero-start mixed prefix cost,
\[
 \limsup_{n\to\infty}\frac{C(n)}{\ln n}=\frac1{\ln4},
 \qquad \liminf_{n\to\infty}\frac{C(n)}{\ln n}=0.
\]
\end{corollary}
\begin{proof}
If $C(n)=k$, Corollary~\ref{cor:first} gives $n\geq g_k$.
As $k\to\infty$, $\ln g_k=(k-1)\ln4+\ln(2/3)+o(1)$,
so $k/\ln g_k\to1/\ln4$; bounded $k$ contributes limit zero.
Equality in the upper limit follows along $n=g_k$.
For the lower limit use $C(2^r)=W(2^r)=1$ and nonnegativity.
\end{proof}

\begin{corollary}\label{cor:types}
A nonempty Thue--Morse prefix is a palindrome exactly when its length
is $4^r$, and an antipalindrome exactly when its length is $2\cdot4^r$,
where $r\geq0$.
\end{corollary}
\begin{proof}
Being allowed is equivalent to $C(n)=1$, hence to $W(n)=1$,
which for positive $n$ means $n=2^e$.
The prefix of length $2^e$ is $\mu^e(0)$.
The reflection identity used in Lemma~\ref{lem:cut} interchanges the
two types at each application of $\mu$, starting with the palindrome $0$.
A nonempty binary word cannot have both types.
The interchange is also \cite[Observation~1(iv)--(v)]{AMP}.
\end{proof}

For the other endpoint we use Frid's published recurrence
\cite[Theorem~4, equations (1)--(4)]{Frid}:
\begin{align}
 P(4m)&=P(m),&P(4m+1)&=P(m)+1,\nonumber\\
 P(4m+2)&=\min\{P(m),P(m+1)\}+2,
 &P(4m+3)&=P(m+1)+1.\label{eq:frid}
\end{align}
Here the word and prefix conventions are the same; Frid uses letters
$a,b$ and the notation $\operatorname{PPL}_t$.

\begin{corollary}\label{cor:sharp}
Let $F(n)=64n+10$. For every $n\geq0$,
\[
 P(F(n))=P(n)+4,\qquad C(F(n))=C(n)+2.
\]
In particular, if $v_0=0$, $z_0=2$, and
$v_{r+1}=F(v_r)$, $z_{r+1}=F(z_r)$, then
\[
 (P(v_r),C(v_r))=(4r,2r),\qquad
 (P(z_r),C(z_r))=(4r+2,2r+1).
\]
\end{corollary}
\begin{proof}
Write $p=P(n)$ and $w=W(n)$. Equations \eqref{eq:frid} and
\eqref{eq:four} successively give the following pairs of adjacent values:
\[
\begin{array}{c|cc|cc}
 j&P(j)&P(j+1)&W(j)&W(j+1)\\\hline
 4n&p&p+1&w&w+1\\
 16n+2&p+2&p+2&w+1&w+2\\
 64n+10&p+4&p+3&w+2&w+3
\end{array}
\]
Use $C=W$ and iterate, starting with $(P(0),C(0))=(0,0)$ and
$(P(2),C(2))=(2,1)$. Positive even costs use $v_r$ with $r\geq1$;
positive odd costs use $z_r$ with $r\geq0$.
\end{proof}

The prefix restriction matters. The shifted factor $t[1:7)=110100$
is an antipalindrome, so its mixed cost is one. Its ordinary cost is
three: $11\mid010\mid0$ is a factorization, and its only proper
nonempty palindromic prefixes are $1$ and $11$, whose complementary
suffixes $10100$ and $0100$ are not palindromes. The whole word is not
a palindrome either. Thus the factor-two inequality is false for
arbitrary shifted factors.

\clearpage
\section{Coverage and verification}
The table separates source assertions, questions and consequences.
Page numbers refer to arXiv:1812.00711v2 \cite{AKMP}.
\begingroup\small
\begin{longtable}{@{}p{0.34\linewidth}p{0.60\linewidth}@{}}
\toprule
Source item&Coverage\\\midrule
Section 6, p. 13; Table 1, p. 14: sharp factor two&
Proved for all zero-start prefixes; Corollary~\ref{cor:sharp} gives
attainment at every positive mixed cost.\\[3pt]
Section 6, p. 13; Table 2, p. 14: unbounded equal costs&
Proved by Corollary~\ref{cor:equal}, including its displayed length family.\\[3pt]
Table 2: first equal-cost occurrences&
Proved for every cost by Corollary~\ref{cor:first}; the underlying signed-weight
record sequence is already OEIS A007583.\\[3pt]
Table 1: first sharp occurrences&
Not covered in general. Its costs 1--5 are finite word checks;
all displayed row values follow from \eqref{eq:frid} and $C=W$.
Corollary~\ref{cor:sharp} asserts attainment, not first occurrence.\\[3pt]
Section 6, Figure 2: mixed prefix growth&
Not posed as a formula; proved in Corollary~\ref{cor:growth}.
For ordinary prefixes, Frid's Corollary~16 already gives
$\limsup P(n)/\ln n=3/(4\ln2)$; prior proof cited.\\[3pt]
Section 4, p. 10: conjectured all-factor Thue--Morse growth;
pp. 9--10: Fibonacci estimates&
Not covered. The maximum over shifted factors is not $P(n)$ or $C(n)$;
the prefix formulas do not determine those constants.\\[3pt]
Section 7(i): powers conjugate to palindrome-image morphisms&
Not covered as a general question. The source's Example~3 already gives
$\mu^2(0)=0110$, $\mu^2(1)=1001$; prior example cited.\\[3pt]
Section 7(ii)--(iii): lower growth methods and Fibonacci factors&
Not covered in their general or Fibonacci scope. For Thue--Morse prefixes,
Frid's Proposition~14 and Corollary~16 give the exact ordinary first
occurrences and upper growth rate; prior proof cited.\\[3pt]
Section 7(iv): uniformly recurrent words with square-root growth&
Not covered; no assertion about such words follows here.\\[3pt]
Section 7(v): morphisms with arbitrarily long antipalindromes&
Not covered as a general classification. Corollary~\ref{cor:types}
classifies the allowed Thue--Morse prefixes. AMP's Theorems~34 and 36
give later uniform and palindromic special cases; prior proof cited.\\[3pt]
Arbitrary-factor version of the factor-two comparison&
Not asserted in the source; fails for $t[1:7)=110100$,
whose ordinary and mixed costs are three and one.\\
\bottomrule
\end{longtable}
\endgroup

\begingroup\small
The pinned
\href{https://github.com/the-omega-institute/trureturing/blob/106c79bbb6571a3dc7c60f4b2146860d56e88fc5/D5/S1/Words/Palindromes/ThueMorseSignedCut.lean}{signed-cut module}
defines \texttt{MixedFactors} as consecutive nonempty allowed factors
and \texttt{combinedLength} as their minimum number.
Its \texttt{allowed\_factor\_signed\_cut} proves \eqref{eq:cut}, and
\texttt{combined\_prefix\_signed\_lower} proves $W(n)\leq C(n)\leq P(n)$.
The pinned
\href{https://github.com/the-omega-institute/trureturing/blob/106c79bbb6571a3dc7c60f4b2146860d56e88fc5/D5/S1/Words/Palindromes/ThueMorseSharpPrefix.lean}{sharp-prefix module}
contains \texttt{prefix\_signed\_upper} for $P(n)\leq2W(n)$ and
\texttt{sharp\_factor\_two\_prefix} for the universal comparison
and its positive witness.
The declarations \texttt{family\_attains} and
\texttt{unbounded\_equal\_cost} in
the pinned \href{https://github.com/the-omega-institute/trureturing/blob/106c79bbb6571a3dc7c60f4b2146860d56e88fc5/D5/S1/Words/Palindromes/ThueMorseAttainment.lean}{attainment module}
give Corollary~\ref{cor:equal} and its unboundedness consequence.
Their binary word is \texttt{factor n 0}, with false representing $0$.
The equality $C=W$, the first-occurrence and growth formulas,
the complete prefix-type classification, Corollary~\ref{cor:sharp},
and the shifted-factor example have written proofs above;
these extensions are distinct from the selected declarations' scope.
\par\endgroup

The accompanying \texttt{verify.py} computes exact ordinary and mixed word
minima and compares them with the signed-weight and Frid recurrences.
It also checks the first equal-cost occurrences and the prefix-type
classification. The finite ranges and outputs are given in the supporting
files; the universal assertions use the written proofs above.

\begin{thebibliography}{9}
\small\setlength{\itemsep}{1pt}
\bibitem{AKMP} P. Ambro\v{z}, O. Kadlec, Z. Mas\'akov\'a and E. Pelantov\'a,
\emph{Palindromic length of words and morphisms in class $\mathcal P$},
Theoretical Computer Science \textbf{780} (2019), 74--83.
\href{https://doi.org/10.1016/j.tcs.2019.02.024}{doi:10.1016/j.tcs.2019.02.024}.
Source text: \href{https://arxiv.org/abs/1812.00711v2}{arXiv:1812.00711v2},
Section 6, pp. 12--14, and Section 7, pp. 14--15.
\bibitem{Frid} A. E. Frid, \emph{Prefix palindromic length of the Thue--Morse
word}, Journal of Integer Sequences \textbf{22} (2019), Article 19.7.8.
\href{https://arxiv.org/abs/1906.09392v2}{arXiv:1906.09392v2},
Lemma 3 and Theorem 4, pp. 2--3.
\bibitem{Li} S. Li, \emph{Palindromic length complexity and a generalization
of Thue--Morse sequences},
\href{https://arxiv.org/abs/1907.12543v1}{arXiv:1907.12543v1} (2019),
Proposition 1.
\bibitem{AMP} P. Ambro\v{z}, Z. Mas\'akov\'a and E. Pelantov\'a,
\emph{Morphisms generating antipalindromic words},
European Journal of Combinatorics \textbf{89} (2020), 103160,
\href{https://doi.org/10.1016/j.ejc.2020.103160}{doi:10.1016/j.ejc.2020.103160};
\href{https://arxiv.org/abs/1906.06174v1}{arXiv:1906.06174v1} (2019),
Observation 1 and Theorems 34 and 36.
\bibitem{OEIS} OEIS Foundation Inc., \emph{The On-Line Encyclopedia of
Integer Sequences}, entries \href{https://oeis.org/A007302}{A007302}
(signed binary weight) and \href{https://oeis.org/A007583}{A007583}
(record positions), accessed 18 November 2026.
\bibitem{KW} S. Kropf and S. Wagner, \emph{$q$-Quasiadditive Functions},
\href{https://arxiv.org/abs/1605.03654v1}{arXiv:1605.03654v1} (2016),
Section~2, p.~3, ``The nonadjacent form and its Hamming weight.''
\bibitem{BBGL} A. Blondin-Mass\'e, S. Brlek, A. Garon and S. Labb\'e,
\emph{Combinatorial properties of $f$-palindromes in the Thue--Morse sequence},
Pure Mathematics and Applications \textbf{19}(2--3) (2008), 39--52.
\href{https://www.slabbe.org/Publications/2008-ThueMorseLacunas.pdf}{Author text},
Section~5, Proposition~6 and Lemmas~2--3, pp.~10--11.
\end{thebibliography}
\end{document}
